% =====================================================================
\chapter{A feladat formális modellje}
\label{ch:formalizmus}
% =====================================================================

A vizsgagenerálás minőségéről csak akkor lehet szigorúan beszélni, ha a „jó vizsga” fogalmát mérhető
predikátumokra bontjuk. Ebben a fejezetben bevezetjük a dokumentum, a darabolás, a vizsgaspecifikáció
és a kérdés formális modelljét, definiáljuk a minőségi predikátumokat és az ezekből képzett mutatókat,
majd a módszer két szerkezeti tulajdonságát bizonyítjuk. A definíciók egyúttal \az{\ref{ch:kiertekeles}}.~fejezet mutatóinak pontos specifikációi is: a kiértékelő keretrendszer
(\code{tests/eval/mimir\_eval}) ezeket valósítja meg.

\section{Dokumentum, mondatok és darabok}

\begin{definicio}[Dokumentum és mondatfelbontás]
A \emph{dokumentum} egy $\Doc\in\Sigma^{*}$ karaktersorozat. A mondatfelbontás a dokumentumot
diszjunkt, egymást követő intervallumokra bontja: $S=(s_1,\dots,s_m)$, ahol minden $s_i$ egy
$[a_i,b_i)$ karakterintervallumhoz tartozik és $b_i\le a_{i+1}$.
\end{definicio}

\begin{definicio}[Darabolás]
A \emph{darabolás} a mondatsorozat egy összefüggő partíciója: $\Chunks=(c_1,\dots,c_K)$, ahol
$c_k=(s_{j_{k-1}+1},\dots,s_{j_k})$ és $0=j_0<j_1<\dots<j_K=m$. A darabolást a vágási pontok
$J=\{j_1,\dots,j_{K-1}\}$ halmaza egyértelműen meghatározza.
\end{definicio}

Legyen $\phi:\Sigma^{*}\to\mathbb{S}^{d-1}$ egységnyi normájú beágyazás, és jelölje
$\langle\cdot,\cdot\rangle$ a skaláris szorzatot, így a koszinusz-hasonlóság egységvektorokra
$\cos(u,v)=\langle u,v\rangle$. A szomszédos mondatok \emph{szemantikus távolsága}
\begin{equation}
\delta_i = 1-\cos\bigl(\phi(s_i),\phi(s_{i+1})\bigr),\qquad i=1,\dots,m-1.
\label{eq:delta}
\end{equation}
A szemantikus darabolás alapgondolata, hogy a témaváltásoknál $\delta_i$ nagy. A pontos vágási
szabályt (küszöb és méretfüggő érzékenység) \az{\ref{sec:darabolas}}.~szakaszban adjuk meg.

\section{Vizsgaspecifikáció és kérdés}

\begin{definicio}[Vizsgaspecifikáció]
A \emph{vizsgaspecifikáció} egy $\sigma=(n,T,\lambda,\ell)$ négyes, ahol $n\in\mathbb{N}^{+}$ a
kérdések száma, $T\subseteq\{\textsf{mcq},\textsf{tf},\textsf{open}\}$ a megengedett típusok
sorozata, $\lambda$ a nehézség (könnyű, közepes vagy nehéz), $\ell$ pedig a nyelv.
\end{definicio}

\begin{definicio}[Kérdés]
Egy \emph{kérdés} a $q=(\theta,t,O,R,b)$ ötös, ahol $\theta\in T$ a típus, $t$ a kérdés törzse,
$O=(o_1,\dots,o_{|O|})$ a válaszlehetőségek sorozata, amelyek közül a $\kappa(o)\in\{0,1\}$ címke jelöli a
helyeseket, $R\subseteq\{1,\dots,K\}$ a kérdés által hivatkozott darabok indexhalmaza, $b$ pedig a
tervezett Bloom-szint. Egyetlen helyes lehetőség esetén ezt \emph{kulcsnak} nevezzük és $o^{\star}$-gal
jelöljük; a többi lehetőség a \emph{disztraktorok} halmaza: $O^{-}=\{o\in O:\kappa(o)=0\}$.
\end{definicio}

A típusonként elvárt opciószámot $m_{\textsf{mcq}}=4$, $m_{\textsf{tf}}=2$, $m_{\textsf{open}}=1$ jelöli.

\begin{definicio}[Formai helyesség]
A $q$ kérdés \emph{formailag helyes}, ha
\begin{equation}
\mathrm{Shape}(q)=\1\Bigl[\,|O|=m_{\theta}\;\wedge\;\textstyle\sum_{o\in O}\kappa(o)=1\Bigr]
\qquad(\theta\in\{\textsf{mcq},\textsf{tf}\}),
\end{equation}
nyitott kérdésnél pedig $\mathrm{Shape}(q)=\1[\,\exists o\in O:\kappa(o)=1 \wedge o\ne\varepsilon\,]$.
Az $\Exam=(q_1,\dots,q_N)$ vizsga \emph{formai megfelelése}
\begin{equation}
F(\Exam)=\1\Bigl[\,N=n\;\wedge\;\forall i:\mathrm{Shape}(q_i)=1\;\wedge\;\forall i:\theta_i\in T\Bigr].
\label{eq:format}
\end{equation}
\end{definicio}

\section{Forrásalapú minőségi predikátumok}

A kérdés tartalmi helyessége a forráshoz viszonyított fogalom. Jelölje $\mathrm{Src}(q)\subseteq\Doc$
azt a szövegrészt, amelyből a kérdés készült (a hivatkozott vagy a kontextusként átadott darabok
konkatenációja), és $\models$ a természetes nyelvi következményrelációt: $X\models y$, ha egy
figyelmes olvasó $X$ alapján $y$-t igaznak fogadja el. A $\models$ reláció nem számítható ki pontosan;
a kiértékelésben egy független bíró (\ref{sec:biro}.~szakasz) közelíti, az oktatói értékelés pedig ezt
a közelítést validálja.

\begin{definicio}[Alátámasztottság]
A $q$ kérdés \emph{alátámasztott}, ha a forrás igazolja a kulcsot:
$G(q)=\1\bigl[\mathrm{Src}(q)\models o^{\star}\text{ mint } t \text{ válasza}\bigr]$.
\end{definicio}

\begin{definicio}[Disztraktor-validitás]
A $q$ kérdés disztraktor-validitása a forrás szerint hamis \emph{és} témába vágó disztraktorok aránya:
\begin{equation}
V(q)=\frac{1}{|O^{-}|}\sum_{o\in O^{-}}\1\bigl[\mathrm{Src}(q)\models\neg o\bigr]\cdot\1\bigl[\mathrm{onTopic}(o,t)\bigr].
\end{equation}
\end{definicio}

\begin{definicio}[Vak megoldhatóság]
Legyen $\pi$ az opciók egyenletes eloszlású véletlen permutációja, és $\mathcal{J}(t,\pi(O),\mathrm{Src})$
egy megoldó választása, amely a kulcs ismerete nélkül, csak a forrás alapján dönt. A kérdés \emph{vakon
megoldható}, ha $A(q)=\1\bigl[\mathcal{J}(t,\pi(O),\mathrm{Src}(q))=o^{\star}\bigr]$.
\end{definicio}

A vak megoldhatóság a kulcs helyességének és a kérdés egyértelműségének együttes, viselkedési mércéje:
ha két opció is védhető, vagy a kulcs hibás, a megoldó várhatóan más opciót választ. A permutáció a
bíró pozíciós torzítását~\cite{wang2024fair} semlegesíti. Ha a megoldó nem ad értelmezhető választ,
a kérdést a vak megoldhatóság átlagából kihagyjuk (nem nullának számítjuk).

\section{Vizsgaszintű mutatók}

\begin{definicio}[Minőségi index]
Egy vizsga \emph{minőségi indexe} a formai megfelelés és a három kérdésszintű ráta átlaga:
\begin{equation}
Q(\Exam)=\frac14\Bigl(F(\Exam)+\overline{G}(\Exam)+\overline{A}(\Exam)+\overline{V}(\Exam)\Bigr),
\qquad \overline{X}(\Exam)=\frac{1}{|\Exam_X|}\sum_{q\in\Exam_X}X(q),
\label{eq:qindex}
\end{equation}
ahol $\Exam_X\subseteq\Exam$ azon kérdések halmaza, amelyekre az $X$ predikátum értelmezett.
\end{definicio}

A $Q$ index szándékosan egyszerű, súlyozatlan átlag: minden komponense $[0,1]$-beli, a nagyobb érték
jobb, és a formai megfelelés súlya egyenlő a tartalmi mutatókéval. Ennek következménye, hogy egyetlen
hibás formájú kérdés egy tízkérdéses vizsgában a $Q$-t $0{,}25$-tel csökkenti -- a pilotban pontosan ez
történt az ellenőrző karok egy-egy vizsgájában (\ref{sec:fallback}.~szakasz).

\begin{definicio}[Lefedettség]
Legyen $\nu(q)=\argmax_{k}\cos\bigl(\phi(t\oplus o^{\star}),\phi(c_k)\bigr)$ a kérdéshez (törzs és kulcs
összefűzéséhez) legközelebbi darab. A vizsga \emph{lefedettsége}
\begin{equation}
\mathrm{Cov}(\Exam)=\frac{\bigl|\{\nu(q):q\in\Exam\}\bigr|}{K}.
\label{eq:cov}
\end{equation}
\end{definicio}

A lefedettség felső korlátja $\min(n,K)/K$; egy tízkérdéses vizsga egy négydarabos dokumentumon elvben
teljes lefedettséget érhet el. A lefedettség nem minőségi, hanem \emph{szerkezeti} mutató: azt méri,
mennyire oszlik el a vizsga a tananyagon, ezért a $Q$ indexbe nem számítjuk be, hanem mellette
közöljük (a P2 probléma közvetlen mérőszáma).

\begin{definicio}[Válaszszivárgás és nem-kérdés törzs]
Jelölje $W(x)$ az $x$ szöveg szótokenjeinek halmazát. A kérdés \emph{szivárogtatja a kulcsot}, ha
\begin{equation}
L(q)=\1\Bigl[\,|W(o^{\star})|\ge 3\;\wedge\;\frac{|W(o^{\star})\cap W(t)|}{|W(o^{\star})|}\ge 0{,}8\Bigr],
\label{eq:leak}
\end{equation}
és a törzs \emph{nem kérdés}, ha feleletválasztós kérdésnél sem kérdőjelre, sem kettőspontra nem végződik:
$N(q)=\1[\,t \text{ utolsó írásjele}\notin\{\text{?},\,\text{:}\}\,]$ (a záró idézőjeleket figyelmen kívül hagyva).
\end{definicio}

A háromszavas alsó korlát a rövid, egyszavas kulcsok (pl.\ „Arabica”) téves pozitív jelzését zárja ki:
egy tulajdonnév természetesen szerepelhet a törzsben. Az $L$ és $N$ predikátumokat a pilot során feltárt
hibamód (\ref{sec:szivargas}.~szakasz) motiválta; azóta mind mutatóként, mind az ellenőrző lánc
szabályaként használjuk őket.

\begin{definicio}[Duplikátumráta]
A vizsga duplikátumrátája azon kérdéspárok aránya, amelyek beágyazásainak koszinusz-hasonlósága
meghaladja a $0{,}9$ küszöböt:
\begin{equation*}
\mathrm{Dup}(\Exam)=\binom{N}{2}^{-1}\sum_{i<j}\1\bigl[\cos(\phi(q_i),\phi(q_j))>0{,}9\bigr].
\end{equation*}
\end{definicio}

\section{A tervezési probléma}

A fenti fogalmakkal a vizsgagenerálás egy korlátos optimalizálási feladatként írható fel. Jelölje
$\Pi$ a csővezetékek (módszerek) terét, $\Exam_\pi(\Doc,\sigma;\omega)$ a $\pi$ csővezeték által
$\omega$ véletlenmaggal előállított vizsgát, $C(\pi)$ a modellhívások számát, $\tau(\pi)$ a futási időt,
$M(\pi)$ pedig a csúcs GPU-memóriát. A cél
\begin{equation}
\max_{\pi\in\Pi}\;\E_{\Doc,\omega}\Bigl[Q\bigl(\Exam_\pi\bigr)\Bigr]
\quad\text{f.h.}\quad
\E\bigl[F(\Exam_\pi)\bigr]=1,\;\;
M(\pi)\le M_{\max},\;\;
\E\bigl[\tau(\pi)\bigr]\le\tau_{\max},
\label{eq:objective}
\end{equation}
másodlagos célként pedig a lefedettség, $\E[\mathrm{Cov}(\Exam_\pi)]$ maximalizálása. A \Mimir{}
esetében $M_{\max}=8$\,GB (egy fogyasztói kártya, amelyen az asztali környezet is fut), $\tau_{\max}$ pedig
egy oktató számára elfogadható várakozási idő (néhány perc egy tízkérdéses vizsgára).

A kísérleti karok (\ref{sec:karok}.~szakasz) a $\Pi$ tér egy-egy pontjai, amelyek egyszerre egy
tényezőben különböznek. \Az{\eqref{eq:objective}} feladat \emph{nem} oldható meg analitikusan, mert $Q$
csak mintavétellel becsülhető; a kiértékelési módszertan éppen ennek a becslésnek a torzítatlanságát és
pontosságát biztosítja.

\section{A vizsgaterv szerkezeti tulajdonságai}
\label{sec:terv-tulajdonsagok}

A tervezett csővezeték két fontos tulajdonsága már a konstrukcióból következik, mérés nélkül is.

\begin{allitas}[Bloom-kvóta]\label{all:kvota}
Legyen $p:\mathcal{L}\to[0,1]$ a nehézséghez rendelt Bloom-eloszlás ($\sum_b p_b=1$), és legyen $c_b$ a
legnagyobb maradék módszerével kapott darabszám: $c_b=\lfloor np_b\rfloor+\1[b\in\mathcal{T}]$, ahol
$\mathcal{T}$ a $\{np_b-\lfloor np_b\rfloor\}$ törtrészek szerint csökkenő sorrendben első
$n-\sum_b\lfloor np_b\rfloor$ szint. Ekkor $\sum_b c_b=n$ és minden $b$-re $|c_b-np_b|<1$.
\end{allitas}

\begin{proof}
A $\lfloor np_b\rfloor$ értékek összege $n-r$, ahol $r=\sum_b(np_b-\lfloor np_b\rfloor)$ egész, mert
$\sum_b np_b=n$; továbbá $0\le r<|\mathcal{L}|$. Pontosan $r$ szint kap $+1$-et, így az összeg $n$.
Ha $b\in\mathcal{T}$, akkor $c_b-np_b=1-(np_b-\lfloor np_b\rfloor)\in(0,1]$, és az egyenlőség csak
egész $np_b$ esetén állna fenn, amikor a törtrész $0$ -- ilyen szint azonban csak akkor kerülhet
$\mathcal{T}$-be, ha minden törtrész $0$, ekkor viszont $r=0$ és $\mathcal{T}=\emptyset$. Ha
$b\notin\mathcal{T}$, akkor $c_b-np_b=-(np_b-\lfloor np_b\rfloor)\in(-1,0]$.
\end{proof}

Az állítás a választási matematikából ismert Hamilton-módszer kvótatulajdonsága~\cite{balinski2001fair}.
Példaként közepes nehézségnél ($p=(0{,}4;\,0{,}4;\,0{,}2)$ a megértés, alkalmazás, elemzés szintekre) és
$n=10$ esetén a kiosztás pontosan $(4,4,2)$; $n=7$ esetén $(2{,}8;\,2{,}8;\,1{,}4)\to(3,3,1)$. A
nehézség így nem kérés, hanem \emph{garantált} szerkezeti tulajdonság.

\begin{allitas}[Tervezett lefedettség]\label{all:lefed}
Legyen $\mathcal{K}=\{\kappa_1,\dots,\kappa_P\}$ a tervező által javasolt fogalmak halmaza, $D(\kappa)\subseteq\{1,\dots,K\}$
a fogalmat alátámasztó darabok halmaza és $U=\bigcup_{\kappa\in\mathcal{K}}D(\kappa)$. A mohó slotkiosztás
(\ref{alg:blueprint}.~algoritmus), amely minden lépésben olyan fogalmat választ, amely legalább egy
még le nem fedett darabot érint, ha van ilyen, $n\le P$ slot után legalább $\min(n,|U|)$ különböző darabot
fed le.
\end{allitas}

\begin{proof}
Jelölje $H_s$ az első $s$ kiválasztott fogalom által lefedett darabok halmazát. Az $(s+1)$-edik lépésben
két eset lehetséges. Ha van olyan megmaradt $\kappa$, amelyre $D(\kappa)\setminus H_s\ne\emptyset$,
akkor az algoritmus ilyet választ, így $|H_{s+1}|\ge|H_s|+1$. Ha nincs ilyen, akkor minden megmaradt fogalom
darabjai már $H_s$-ben vannak, és a már kiválasztottaké definíció szerint is, így $H_s\supseteq U$, azaz
$|H_s|=|U|$. Indukcióval $|H_n|\ge\min(n,|U|)$.
\end{proof}

Megjegyezzük, hogy a klasszikus maximális lefedési feladatra a határhaszon szerinti mohó algoritmus
$(1-1/e)$-approximációt ad~\cite{nemhauser1978}. A \Mimir{} kiosztása ettől szándékosan eltér: a
választás lexikografikus kulcsa (\emph{új darabot érint-e}, \emph{fontosság}), vagyis az újdonság csak
bináris feltétel, és azon belül a tervező által becsült fontosság dönt. Ezzel az approximációs garanciát a
tartalmi relevanciáért cseréljük, \az{\ref{all:lefed}}.~állítás alsó korlátja azonban megmarad. A tervezett
lefedettség (a kontextusként felhasznált darabok halmaza) és \az{\eqref{eq:cov}} szerint mért lefedettség
nem azonos: az utóbbi azt méri, miről szól ténylegesen a kész kérdés.

\section{A modellhívások számának korlátai}
\label{sec:hivasszam}

A tervezett csővezeték költségét elsősorban a nyelvimodell-hívások száma határozza meg. Legyen $r$ a
slotonkénti újrapróbálkozások maximális száma, és tegyük fel, hogy minden hívás legfeljebb egyszer
ismétlődik értelmezhetetlen JSON miatt. Egy \emph{kísérlet} legfeljebb három sikeres hívásból áll
(generálás, alátámasztottság-ellenőrzés, vak megoldás), egy slot legfeljebb $r+1$ kísérletből, sikertelenség
esetén pedig egy tartalékfogalommal újabb legfeljebb $r+1$ kísérletből és egy javító hívásból. Így az
ellenőrzött csővezeték hívásszámára
\begin{equation}
1+3n\;\le\;C_{\textsf{E2}}\;\le\;2+n\bigl(2\cdot 2\cdot 3\,(r+1)+2\bigr),
\label{eq:callbounds}
\end{equation}
ahol a felső korlátban a tényezők rendre a tartalékfogalmat, a JSON-ismétlést és a kísérletenkénti
három hívást, a $+2$ a javító hívást, a konstans $2$ pedig a (legfeljebb egyszer ismételt) tervezőhívást
jelöli; az alsó korlát az az eset, amikor minden kérdés elsőre átmegy (egy tervezőhívás, majd
slotonként három hívás). Ha egy kísérlet $p$ valószínűséggel megy át, függetlenül a többitől, akkor a
slotonkénti kísérletek várható száma a csonkolt geometriai eloszlás várható értéke:
\begin{equation}
\E[a]=\sum_{j=0}^{r}(1-p)^{j}=\frac{1-(1-p)^{r+1}}{p}.
\label{eq:expattempts}
\end{equation}
Az egyhívásos karok (\arm{E0}, \arm{E5a--c}, \arm{B-doc}) hívásszáma 1, az ellenőrzés
nélküli \arm{E1} kar alsó korlátja $1+n$, az \arm{E2f} karé $1+2n$, a gráfalapú \arm{E3} karé pedig
$G+3n$, ahol $G=\lceil|\Doc|/3000\rceil$ a gráfkötegek száma. A pilotban ($n=10$, $r=2$) az \arm{E2} kar átlagosan 75 hívást igényelt, ami slotonként 7,4 hívásnak
felel meg -- ez messze az alsó korlát (3) felett van, és jelzi, hogy a kis modell ellenőrző lánca a
kérdések jelentős részét elutasította (\ref{sec:koltseg}.~szakasz). A naiv alapvonal ezzel szemben
$C_{\textsf{E0}}=1$ hívással dolgozik; a két módszer közötti költségkülönbség tehát két nagyságrend,
amelyet a minőségi nyereségnek kell igazolnia.
